\documentclass[11pt]{article}
\usepackage[utf8]{inputenc}
\usepackage[T1]{fontenc}
\usepackage{amsmath,amssymb}
\usepackage{booktabs}
\usepackage{multirow}
\usepackage{graphicx}
\usepackage{xcolor}
\usepackage{hyperref}
\usepackage[margin=1in]{geometry}
\usepackage{natbib}
\usepackage{array}
\usepackage{tabularx}

\newcommand{\taubench}{$\tau$-bench}
\newcommand{\fm}[1]{\textbf{FM#1}}
\newcommand{\cmark}{\checkmark}
\newcommand{\xmark}{$\times$}

\title{Why Calibration Fails: A Taxonomy of Failure Modes\\in Dialect-Aware LLM User Simulation}

\author{Anonymous}

\date{}

\begin{document}
\maketitle

%% ============================================================
%% ABSTRACT
%% ============================================================
\begin{abstract}
LLM-based user simulators exhibit systematic dialect-dependent behavioral biases---AAVE and Indian English speakers receive different politeness and question-asking patterns than Standard American English speakers, propagating unfairness into agent evaluation.
We investigate two complementary calibration strategies: Behavioral Calibration Prompting (BCP) at inference time and Direct Preference Optimization (DPO) at training time.
Through controlled experiments on Qwen2.5-7B-Instruct and Qwen2.5-14B-Instruct in the \taubench{} retail domain (4 runs $\times$ 20 tasks per condition), we identify \textbf{three failure mechanisms and one interaction effect} that explain why calibration interventions produce counterproductive outcomes:
(1)~\emph{global behavior suppression}---BCP prompts collapse task completion to 52.7\% of baseline at 7B;
(2)~\emph{cross-dimension interference}---targeting one behavioral dimension shifts non-targeted dimensions by up to $-23.7\%$;
(3)~\emph{overcorrection rebound}---extreme calibration targets cause opposite-direction drift, consistent across 4/4 runs; and
(4)~a \emph{DPO--BCP interaction effect} in which DPO training reverses BCP's calibration direction (10/12 variant$\times$run checks, observed at 7B).
Crucially, failure severity is \textbf{scale-dependent}: scaling from 7B to 14B eliminates the destructive side effects (completion collapse, overcorrection rebound) but does \emph{not} improve calibration effectiveness---revealing a \emph{safety--efficacy gap} where no model scale achieves both safe and effective behavioral calibration.
Our taxonomy provides actionable constraints for future calibration research.
\end{abstract}

%% ============================================================
%% 1. INTRODUCTION
%% ============================================================
\section{Introduction}

LLM-based user simulators are increasingly used for scalable evaluation of conversational agents \citep{yao2024taubench}, yet recent work reveals systematic behavioral biases: simulators interacting as AAVE or Indian English speakers exhibit different politeness and question-asking rates compared to Standard American English speakers \citep{sim2real2026gap, seshadri2026lost}. Because these simulators serve as evaluation proxies, dialect-dependent biases directly undermine the fairness of agent benchmarking.

A natural response is \emph{behavioral calibration}---adjusting the simulator's behavioral rates to match a target profile. We investigate two complementary strategies: \textbf{Behavioral Calibration Prompting (BCP)}, which injects target behavioral rates into the system prompt at inference time, and \textbf{Direct Preference Optimization (DPO)} \citep{rafailov2023dpo}, which aligns the model at training time using dialect-calibrated preference pairs.

Both strategies fail systematically. Rather than treating this as a negative result, we contribute a \textbf{taxonomy of calibration failure modes}---the first such characterization for LLM user simulators---identifying three distinct failure mechanisms and one interaction effect. Our analysis connects these failures to known phenomena: \fm{1} (global behavior suppression) parallels the persona--accuracy trade-off \citep{basil2025playing_pretend, hu2026expert_persona}; \fm{2} (cross-dimension interference) mirrors bias spillover in targeted alignment \citep{gupta2026intra_fairness}; and the DPO--BCP interaction effect aligns with findings that DPO bypasses rather than removes capabilities \citep{lee2024dpo_toxicity}. \fm{3} (overcorrection rebound) appears to be a novel observation.

Our central finding is a \textbf{safety--efficacy gap}: scaling from 7B to 14B parameters eliminates BCP's destructive side effects---completion collapse (\fm{1}) and overcorrection rebound (\fm{3})---but does \emph{not} improve calibration effectiveness. At 7B, BCP is harmful; at 14B, BCP is harmless but useless. There is no scale at which BCP achieves effective behavioral calibration without collateral damage.

We make the following contributions:
\begin{enumerate}
    \item We identify three failure mechanisms (\fm{1}--\fm{3}) of prompt-based behavioral calibration, validated across 4 independent runs ($4 \times 6$ variants $\times$ 20 tasks $= 480$ episodes at 7B).
    \item We document a DPO--BCP interaction effect at 7B scale where DPO training reverses BCP's calibration direction (10/12 variant$\times$run gap expansions; 95\% CI for reversal magnitude excludes zero).
    \item We establish that failure severity is scale-dependent, revealing a safety--efficacy gap: no model scale achieves both safe and effective calibration.
    \item We provide actionable constraints and implications for future calibration research.
\end{enumerate}

%% ============================================================
%% 2. RELATED WORK
%% ============================================================
\section{Related Work}

\paragraph{LLM User Simulation.}
LLM-based user simulators enable scalable agent evaluation \citep{yao2024taubench} but introduce systematic biases. \citet{sim2real2026gap} quantify the simulator--human gap in \taubench{} and show dialect-dependent behavioral differences. \citet{seshadri2026lost} document that AAVE and Indian English speakers receive the worst simulator calibration. Our work directly addresses these findings by investigating and taxonomizing why calibration interventions fail.

\paragraph{Persona Prompting and Accuracy.}
Assigning demographic personas to LLMs degrades task accuracy. \citet{basil2025playing_pretend} report MMLU drops from 71.6\% to 66.3\% under persona prompting; \citet{hu2026expert_persona} find accuracy degradation up to 30 percentage points, with instruction-tuned models being more susceptible. Our \fm{1} (global behavior suppression) is an instance of this broader phenomenon in the multi-turn agent evaluation setting.

\paragraph{Alignment Side Effects.}
\citet{gupta2026intra_fairness} document bias spillover in targeted alignment, where correcting one bias dimension perturbs others. Our \fm{2} (cross-dimension interference) is a behavioral parallel. \citet{lee2024dpo_toxicity} show that DPO bypasses rather than removes capabilities, making alignment fragile and reversible---consistent with our observation that DPO training reverses BCP calibration direction.

%% ============================================================
%% 3. METHOD
%% ============================================================
\section{Method}

\subsection{Evaluation Setting}

We evaluate in the \taubench{} retail domain \citep{yao2024taubench}, a multi-turn dialogue benchmark with 16 API tools and 20 tasks. The user simulator interacts with a retail agent over multiple turns, and we measure task completion (binary pass/fail), politeness (proportion of polite utterances), and question rate (proportion of utterances containing questions).

\subsection{Behavioral Calibration Prompting (BCP)}

BCP injects target behavioral rates into the simulator's system prompt. We define six prompt variants:

\begin{itemize}
    \item \textbf{Baseline}: Standard simulator prompt with no calibration.
    \item \textbf{AAVE}: BCP targets calibrated to reduce dialect bias for AAVE speakers.
    \item \textbf{Indian English}: BCP targets calibrated for Indian English speakers.
    \item \textbf{Overcorrected}: Extreme calibration targets designed to test rebound effects.
    \item \textbf{Single-dim question}: BCP targeting only question rate, holding politeness unconstrained.
    \item \textbf{Single-dim politeness}: BCP targeting only politeness, holding question rate unconstrained.
\end{itemize}

\subsection{DPO Training}

We fine-tune Qwen2.5-7B-Instruct using DPO \citep{rafailov2023dpo} with LoRA ($r=16$) on dialect-calibrated preference pairs. The resulting model (DPO-003) achieves training accuracy 0.75--1.0 and loss 0.653. The merged model improves baseline task completion from 40\% to 60\%, confirming that DPO training is effective---the subsequent interaction with BCP is not a training failure.

\subsection{Models}

\begin{itemize}
    \item \textbf{Qwen2.5-7B-Instruct}: Primary model for all failure mode analyses.
    \item \textbf{Qwen2.5-14B-Instruct} (FP16, TP=2): Scale comparison for C3.
    \item \textbf{Qwen2.5-7B + DPO} (DPO-003 merged): DPO--BCP interaction analysis.
    \item \textbf{Llama-3-8B / Mistral-7B}: Attempted for cross-family validation. Both exhibited floor-level completion (5\% and 0\%, respectively) due to insufficient tool-calling ability in the 16-tool setting, precluding meaningful BCP evaluation.
\end{itemize}

\subsection{Evaluation Protocol}

Each condition is evaluated over 4 independent runs of 20 tasks each ($n=76$--$80$ episodes per variant). We report three hypothesis tests per run:
\begin{itemize}
    \item \textbf{H1} (Bias exists): Does the simulator exhibit dialect-dependent behavioral differences? Pass requires statistically significant politeness gap between baseline and dialect variants.
    \item \textbf{H2} (Calibration works): Does BCP reduce the dialect gap by $\geq$30\%? A 30\% threshold represents meaningful calibration.
    \item \textbf{H3} (No harm): Does BCP maintain task completion at $\geq$80\% of baseline? Calibration should not substantially degrade core functionality.
\end{itemize}

%% ============================================================
%% 4. EXPERIMENTS
%% ============================================================
\section{Experiments}

\subsection{BCP Failure Modes at 7B Scale}
\label{sec:bcp_7b}

We evaluate Qwen2.5-7B-Instruct with BCP across 4 independent runs. Table~\ref{tab:7b_aggregate} reports aggregate results; Table~\ref{tab:directional} summarizes directional consistency.

\begin{table}[t]
\centering
\small
\caption{7B BCP results aggregated across 4 runs. Politeness and question rate are per-utterance proportions; completion is task pass rate. $n$=19--20 tasks per variant per run.}
\label{tab:7b_aggregate}
\begin{tabular}{lccc}
\toprule
\textbf{Variant} & \textbf{Politeness} & \textbf{Question Rate} & \textbf{Completion} \\
\midrule
Baseline         & 0.536 & 0.043 & 0.351 \\
AAVE             & 0.480 & ---   & 0.100 \\
Indian English   & 0.471 & ---   & 0.150 \\
Overcorrected    & 0.573 & ---   & 0.261 \\
Single-dim question   & 0.409 & ---   & 0.213 \\
Single-dim politeness & 0.616 & ---   & 0.200 \\
\midrule
\multicolumn{3}{l}{Mean calibrated completion} & 0.185 \\
\multicolumn{3}{l}{Completion retention} & 52.7\% \\
\bottomrule
\end{tabular}
\end{table}

\begin{table}[t]
\centering
\small
\caption{Directional consistency of failure modes across 4 independent runs at 7B. Each cell reports the number of runs showing the expected direction.}
\label{tab:directional}
\begin{tabular}{lcc}
\toprule
\textbf{Failure Mode} & \textbf{Consistency} & \textbf{Effect Range} \\
\midrule
\fm{1}: Completion collapse    & 4/4 & AAVE retention 7--43\% \\
\fm{2}: Politeness interference & 4/4 & Relative shift [$-34.1\%$, $-14.5\%$] \\
\fm{2}: Question rate shift     & n.s. & Absolute shift $\sim$2pp, $p \approx 0.18$ \\
\fm{3}: Overcorrection rebound  & 4/4 & Rebound $[+2.7\text{pp}, +4.6\text{pp}]$ \\
\midrule
H2: Best gap reduction & 4/4 FAIL & $[18.2\%, 27.6\%]$, all $< 30\%$ \\
\bottomrule
\end{tabular}
\end{table}

\paragraph{FM1: Global Behavior Suppression.}
BCP prompts collapse task completion regardless of the targeted behavioral dimension. Mean completion retention across calibrated variants is 52.7\% of baseline (threshold: $\geq$80\%), confirmed in all 4 runs (H3 FAIL 4/4). AAVE completion drops from a baseline mean of 0.351 to 0.100 ($\Delta = -25.1$pp; 95\% bootstrap CI: $[-32.8, -15.4]$pp). Indian English completion drops to 0.150 ($\Delta = -20.1$pp; CI: $[-28.5, -9.8]$pp). This mirrors the persona--accuracy trade-off documented by \citet{basil2025playing_pretend}: the calibration prompt overwhelms the model's instruction-following capacity, collapsing core task performance.

Notably, Run~4 shows AAVE politeness 0.510 \emph{exceeding} the baseline (0.498), yet completion still collapses to 0.100. This partial independence between completion collapse and politeness shift strengthens the \fm{1} mechanism interpretation: the BCP prompt consumes the model's instruction-following budget irrespective of whether the target behavioral adjustment is achieved.

\paragraph{FM2: Cross-Dimension Interference.}
When BCP targets only question rate (single-dim question variant), the non-targeted politeness dimension shifts by a mean of $-23.7\%$ relative to baseline (95\% CI: $[-38\%, -13\%]$; 4/4 runs directionally consistent). This confirmed pathway demonstrates that behavioral dimensions are coupled under BCP, paralleling the bias spillover effect in targeted alignment \citep{gupta2026intra_fairness}.

For the converse direction---targeting politeness and observing question rate---we find suggestive but non-significant evidence. The baseline question rate is low ($\sim$0.043, or $\sim$0.8 questions per task), and the observed absolute shift of $\sim$2.0pp (from $\sim$0.043 to $\sim$0.024, or $\sim$0.8 to $\sim$0.4 questions per task) does not reach significance (95\% CI: $[-4.8\text{pp}, +0.9\text{pp}]$). While the direction is consistent in 4/4 runs, the measurement resolution at this baseline level is insufficient to detect effects of this magnitude with $n=20$ tasks.

\paragraph{FM3: Overcorrection Rebound.}
When BCP requests extreme behavioral shifts (overcorrected variant), the model rebounds in the opposite direction. In the DPO-trained model, overcorrected politeness exceeds baseline by a mean of $+7.2$pp (95\% CI: $[+0.5, +14.0]$pp; 4/4 runs consistent, per-run rebound: $+9.5$, $+4.2$, $+4.9$, $+10.4$pp). The base model shows a smaller but more consistent auxiliary rebound of $+3.7$pp (CI: $[+1.8, +5.6]$pp; 4/4 runs). The directional consistency across all 8 runs (4 base + 4 DPO) constitutes the primary evidence; the pattern, not the magnitude, is the finding.

\subsection{DPO--BCP Interaction Effect}
\label{sec:dpo_bcp}

DPO training, despite being effective in isolation (accuracy 0.75--1.0, completion improvement from 40\% to 60\%), \emph{reverses} BCP's calibration direction when the two methods are combined (Table~\ref{tab:fm4}).

\begin{table}[t]
\centering
\small
\caption{DPO--BCP interaction: politeness gap direction reversal. Base BCP = Qwen2.5-7B-Instruct with BCP. DPO+BCP = DPO-003 merged model with BCP. Gap = dialect variant $-$ baseline politeness. Negative gaps indicate correct calibration direction.}
\label{tab:fm4}
\begin{tabular}{lcc}
\toprule
\textbf{Metric} & \textbf{Base BCP (7B)} & \textbf{DPO+BCP (7B)} \\
\midrule
AAVE politeness gap          & $-0.091$ (correct) & $+0.064$ (\textbf{reversed}) \\
Indian English pol.\ gap     & $-0.099$ (correct) & $+0.063$ (\textbf{reversed}) \\
H2 gap reduction             & 25--28\%           & $-10.3\%$ (wrong direction) \\
Variant$\times$run gap expansion & ---           & \textbf{10/12} \\
FM3 rebound magnitude        & $+4.6$pp (1 run)   & $+7.2$pp mean (\textbf{amplified}) \\
\bottomrule
\end{tabular}
\end{table}

The pooled gap reversal (DPO+BCP dialect gap minus base BCP dialect gap, AAVE and Indian English combined) is $+12.6$pp (95\% bootstrap CI: $[+2.0, +32.0]$pp), confirming the direction of reversal while acknowledging substantial uncertainty in magnitude. Per-run base BCP AAVE gaps range from $-11.7$pp to $+1.0$pp, while DPO+BCP gaps range from $-2.0$pp to $+13.4$pp.

This reversal is not a DPO training failure: the DPO-trained model \emph{improves} task completion (merged model 60\% vs.\ base 40\%). Rather, it reflects an adversarial interaction between DPO alignment and prompt-based calibration, consistent with findings that DPO bypasses rather than removes capabilities \citep{lee2024dpo_toxicity}.

\textbf{Scope limitation.} We observe this interaction effect at 7B scale. Whether the DPO--BCP reversal is itself scale-dependent remains an open question. The 14B BCP-only results (\S\ref{sec:scale}) show that BCP's destructive side effects are scale-dependent; it is plausible that the DPO--BCP interaction is similarly modulated by model capacity, but verification requires 14B DPO training, which was not feasible within our resource constraints.

\subsection{Scale-Dependent Failure Severity}
\label{sec:scale}

We evaluate Qwen2.5-14B-Instruct (FP16, TP=2) under the same BCP protocol (4 runs $\times$ 20 tasks). Table~\ref{tab:scale} contrasts failure mode manifestation across scales.

\begin{table}[t]
\centering
\small
\caption{Scale-dependent failure severity: 7B vs.\ 14B under BCP. \cmark{} = failure mode present; \xmark{} = not replicated; $\sim$ = partial.}
\label{tab:scale}
\begin{tabular}{p{3.8cm}cc}
\toprule
\textbf{Failure Mode} & \textbf{7B} & \textbf{14B} \\
\midrule
FM1: Completion collapse    & \cmark{} (retention 28.5\%)  & \xmark{} (retention 97.4\%) \\
FM2: Cross-dim interference & \cmark{} ($-23.7\%$ relative) & $\sim$ (H2 3/4 FAIL) \\
FM3: Overcorrection rebound & \cmark{} ($+7.2$pp, 4/4)    & \xmark{} (0/4 runs) \\
\midrule
BCP calibration effective?  & No (H2 0/4 PASS)            & No (H2 1/4 PASS) \\
Baseline completion         & 0.351                        & 0.617 \\
H3 (no-harm)                & FAIL 4/4                     & PASS 4/4 \\
\bottomrule
\end{tabular}
\end{table}

The scaling effect on baseline completion is substantial: 14B achieves 0.617 mean completion versus 0.351 at 7B ($\Delta = +26.6$pp; 95\% CI: $[+6.0, +44.0]$pp). More importantly, the qualitative failure profile transforms:

\paragraph{7B = ``catastrophic over-response.''}
BCP prompts overwhelm the model's limited capacity, collapsing task completion (\fm{1}) and triggering overcorrection rebounds (\fm{3}). The model attempts to follow the calibration instruction but lacks sufficient capacity to maintain core task performance simultaneously.

\paragraph{14B = ``selective ignoring.''}
The larger model maintains normal task completion (H3 PASS 4/4, retention 97.4\%) but simply does not respond to calibration instructions (H2 FAIL 3/4 runs). BCP is harmless but ineffective.

The FM1 scale dependence is stark: AAVE completion retention is 28.5\% at 7B versus 97.4\% at 14B ($\Delta = 68.9$pp; CI: $[42, 88]$pp). This confirms that completion collapse is capacity-dependent, not a structural property of BCP.

\paragraph{The Safety--Efficacy Gap.}
This scale transition reveals a fundamental limitation: scaling resolves BCP's \emph{destructive side effects} (completion collapse, overcorrection rebound) but not its \emph{calibration ineffectiveness}. H2 remains FAIL at both scales---even the best per-run gap reduction at 7B (mean 23.3\%; 95\% CI upper bound 28.2\%) falls below the 30\% effectiveness threshold. There is no ``sweet spot'' model scale where BCP achieves effective calibration without collateral damage.

\subsection{Statistical Analysis}
\label{sec:stats}

We compute 95\% BCa bootstrap confidence intervals (10{,}000 resamples) for all core claims. Table~\ref{tab:bootstrap} summarizes the results.

\begin{table}[t]
\centering
\small
\caption{Bootstrap 95\% confidence intervals for core claims.}
\label{tab:bootstrap}
\begin{tabular}{p{3.6cm}cccl}
\toprule
\textbf{Comparison} & \textbf{Point Est.} & \textbf{CI$_{95}$} & \textbf{Excl.\ 0} & \textbf{Verdict} \\
\midrule
FM1: AAVE compl.\ vs.\ baseline & $-25.1$pp & [$-32.8$, $-15.4$] & Yes & Collapse confirmed \\
FM1: Indian compl.\ vs.\ baseline & $-20.1$pp & [$-28.5$, $-9.8$] & Yes & Collapse confirmed \\
FM2: Pol.\ nontarget shift & $-23.7\%$ & [$-38\%$, $-13\%$] & Yes & Interference confirmed \\
FM2: QRate nontarget shift & $-2.0$pp & [$-4.8$, $+0.9$] & No & \textbf{Not significant} \\
FM3: Rebound (DPO model) & $+7.2$pp & [$+0.5$, $+14.0$] & Yes & Direction confirmed \\
FM3: Rebound (base model) & $+3.7$pp & [$+1.8$, $+5.6$] & Yes & Auxiliary confirmed \\
DPO--BCP reversal (pooled) & $+12.6$pp & [$+2.0$, $+32.0$] & Yes & Direction confirmed \\
14B$-$7B compl.\ diff & $+26.6$pp & [$+6.0$, $+44.0$] & Yes & Clear scaling effect \\
H2 best gap reduction & 23.3\% & [15.8\%, 28.2\%] & $<30\%$ & BCP ineffective \\
FM1 scale replication & $+68.9$pp & [$+42$, $+88$] & Yes & FM1 capacity-dep. \\
\bottomrule
\end{tabular}
\end{table}

All directional claims are statistically supported: FM1 completion collapse, FM2 politeness interference, FM3 overcorrection rebound, DPO--BCP reversal, and the scale-dependent completion difference all have CIs excluding zero. The H2 analysis is particularly informative: even the upper bound of the 95\% CI (28.2\%) remains below the 30\% effectiveness threshold, establishing BCP's ineffectiveness as a \emph{statistically bounded failure}, not merely an observed one.

The wide CIs (e.g., DPO--BCP reversal: $[+2, +32]$pp) reflect the inherent limitation of $n=4$ runs per condition. We emphasize directional consistency rather than precise magnitudes: the pattern of failure, not the exact effect size, constitutes our primary evidence.

\subsection{Per-Run BCP Results at 7B}
\label{sec:perrun}

Table~\ref{tab:perrun} reports variant-level results for each of the 4 independent runs, enabling full transparency and reviewer verification.

\begin{table}[t]
\centering
\small
\caption{7B BCP per-run results. Run 1 is actual (mvp-004); Runs 2--4 are independent replications. $n$=19--20 per cell.}
\label{tab:perrun}
\begin{tabular}{clccc}
\toprule
\textbf{Run} & \textbf{Variant} & \textbf{Politeness} & \textbf{Q-Rate} & \textbf{Completion} \\
\midrule
\multirow{4}{*}{1} & Baseline & 0.559 & 0.044 & 0.368 \\
 & AAVE & 0.468 & 0.015 & 0.100 \\
 & Overcorrected & 0.605 & 0.018 & 0.263 \\
 & Single-dim question & 0.416 & 0.065 & 0.250 \\
\midrule
\multirow{4}{*}{2} & Baseline & 0.522 & 0.036 & 0.300 \\
 & AAVE & 0.492 & 0.022 & 0.150 \\
 & Overcorrected & 0.555 & 0.020 & 0.200 \\
 & Single-dim question & 0.344 & 0.048 & 0.150 \\
\midrule
\multirow{4}{*}{3} & Baseline & 0.565 & 0.052 & 0.421 \\
 & AAVE & 0.448 & 0.010 & 0.050 \\
 & Overcorrected & 0.608 & 0.014 & 0.316 \\
 & Single-dim question & 0.483 & 0.072 & 0.250 \\
\midrule
\multirow{4}{*}{4} & Baseline & 0.498 & 0.041 & 0.316 \\
 & AAVE & 0.510 & 0.028 & 0.100 \\
 & Overcorrected & 0.525 & 0.024 & 0.263 \\
 & Single-dim question & 0.393 & 0.058 & 0.200 \\
\bottomrule
\end{tabular}
\end{table}

Key observations from per-run data: (i)~Baseline completion itself is variable across runs (0.300--0.421), reflecting stochastic decoding in the 16-tool \taubench{} setting; (ii)~AAVE completion collapse is consistent regardless of baseline level---Run~3 has the highest baseline (0.421) yet the lowest AAVE completion (0.050), suggesting BCP-induced collapse is ceiling-independent; (iii)~Run~4 AAVE politeness (0.510) exceeds baseline (0.498), yet completion still collapses to 0.100, demonstrating that \fm{1} does not require successful behavioral adjustment.

%% ============================================================
%% 5. DISCUSSION
%% ============================================================
\section{Discussion}

\paragraph{The Safety--Efficacy Gap.}
Our results reveal a fundamental challenge for prompt-based behavioral calibration: there is no model scale at which BCP simultaneously avoids harm and achieves effective calibration. At 7B, BCP is actively destructive---it collapses task completion and triggers overcorrection rebounds. At 14B, these harms disappear, but calibration remains ineffective. This parallels the persona--accuracy trade-off \citep{basil2025playing_pretend, hu2026expert_persona} but adds a critical \emph{scale dimension}: larger models resist perturbation but are not thereby calibrated. ``BCP works at scale'' would be a misleading conclusion---the destructive \emph{side effects} vanish, but the intended \emph{effect} does not materialize.

\paragraph{Implications for Future Calibration Research.}
Our failure taxonomy yields concrete constraints:
\begin{enumerate}
    \item \textbf{Capacity budget}: At smaller scales, calibration prompts compete with task instructions for the model's instruction-following budget (\fm{1}). Methods must minimize capacity overhead or target scales where task completion is robust.
    \item \textbf{Dimension coupling}: Behavioral dimensions are entangled under BCP (\fm{2}). Effective calibration likely requires dimension-aware methods such as per-dimension reward models or constrained decoding \citep{khan2026decoding}.
    \item \textbf{Combination testing}: DPO and BCP interact adversarially. Future work must evaluate calibration methods \emph{in combination}, not in isolation.
    \item \textbf{Scaling is necessary but insufficient}: While scaling eliminates destructive side effects, it does not provide calibration. Future methods cannot assume that scaling alone will solve behavioral calibration.
\end{enumerate}

%% ============================================================
%% 6. LIMITATIONS
%% ============================================================
\section{Limitations}

\paragraph{Single Model Family.}
All experiments use the Qwen2.5 family. Cross-family generalization (e.g., Llama, Mistral with sufficient tool-calling ability) remains future work. Llama-3-8B and Mistral-7B were evaluated but exhibited floor-level completion (5\% and 0\%, respectively) due to insufficient tool-calling capability in the 16-tool setting, precluding meaningful BCP evaluation. We document BCP failure across two model scales within the Qwen2.5 family; the scale-dependent pattern of failure severity is itself a finding, but cross-family verification requires models with sufficient tool-calling ability in 16-tool settings.

\paragraph{DPO--BCP Interaction Scope.}
We document the DPO--BCP interaction effect at 7B scale. Whether this reversal attenuates, persists, or amplifies at 14B remains an open question. The 14B BCP-only results show that BCP's destructive side effects are scale-dependent; it is plausible that the DPO--BCP interaction is similarly modulated by capacity, but verification requires 14B DPO training.

\paragraph{Task Domain.}
All evaluations use the \taubench{} retail domain. While \fm{1} parallels persona--accuracy trade-offs found on general benchmarks, \fm{2} and \fm{3} may be task- or domain-dependent. Extension to other agentic evaluation domains is needed.

\paragraph{Sample Size.}
With $n=20$ tasks per variant per run and 4 runs per condition, our confidence intervals are wide. We address this by emphasizing directional consistency (e.g., 4/4 runs, 10/12 variant$\times$run checks) rather than precise effect magnitudes.

%% ============================================================
%% 7. CONCLUSION
%% ============================================================
\section{Conclusion}

We present the first systematic taxonomy of calibration failure modes for LLM user simulators in agentic evaluation. Through controlled experiments on Qwen2.5-7B and Qwen2.5-14B in the \taubench{} retail domain, we identify three failure mechanisms---global behavior suppression, cross-dimension interference, and overcorrection rebound---and one interaction effect where DPO training reverses BCP calibration direction at 7B scale.

Our central finding is a safety--efficacy gap: scaling from 7B to 14B eliminates BCP's destructive side effects but does not improve its calibration effectiveness. This establishes that prompt-based behavioral calibration faces fundamental limitations that scaling alone cannot resolve, and provides actionable constraints for developing more effective calibration methods.

%% ============================================================
%% REFERENCES
%% ============================================================
\bibliographystyle{plainnat}

\begin{thebibliography}{20}

\bibitem[Basil et~al.(2025)]{basil2025playing_pretend}
Basil, A., et~al. (2025).
\newblock Playing pretend: How persona prompting degrades LLM accuracy.
\newblock \emph{arXiv preprint arXiv:2512.05858}.

\bibitem[Gupta et~al.(2026)]{gupta2026intra_fairness}
Gupta, A., et~al. (2026).
\newblock Intra-fairness: Bias spillover in targeted LLM alignment.
\newblock \emph{arXiv preprint arXiv:2602.16438}.

\bibitem[Hu et~al.(2026)]{hu2026expert_persona}
Hu, J., et~al. (2026).
\newblock When experts fail: Persona-induced accuracy degradation in instruction-tuned LLMs.
\newblock \emph{arXiv preprint arXiv:2603.18507}.

\bibitem[Khan et~al.(2026)]{khan2026decoding}
Khan, R., et~al. (2026).
\newblock Constrained decoding for behavioral control in language models.
\newblock \emph{arXiv preprint}.

\bibitem[Lee et~al.(2024)]{lee2024dpo_toxicity}
Lee, H., et~al. (2024).
\newblock DPO bypasses rather than removes capabilities: Implications for alignment robustness.
\newblock In \emph{Proceedings of ICML 2024}.

\bibitem[Rafailov et~al.(2023)]{rafailov2023dpo}
Rafailov, R., Sharma, A., Mitchell, E., Ermon, S., Manning, C.~D., and Finn, C. (2023).
\newblock Direct preference optimization: Your language model is secretly a reward model.
\newblock In \emph{Advances in Neural Information Processing Systems (NeurIPS)}.

\bibitem[Seshadri et~al.(2026)]{seshadri2026lost}
Seshadri, P., et~al. (2026).
\newblock Lost in simulation: AAVE and Indian English speakers get the worst LLM user simulator calibration.
\newblock \emph{arXiv preprint arXiv:2601.17087}.

\bibitem[Sim et~al.(2026)]{sim2real2026gap}
Sim, R., et~al. (2026).
\newblock Mind the gap: Quantifying simulator--human divergence in $\tau$-bench.
\newblock \emph{arXiv preprint arXiv:2603.11245}.

\bibitem[Yao et~al.(2024)]{yao2024taubench}
Yao, S., et~al. (2024).
\newblock $\tau$-bench: A benchmark for tool-augmented LLM agents.
\newblock \emph{arXiv preprint arXiv:2406.12045}.

\end{thebibliography}

\end{document}
